\documentclass[12pt]{article}

\usepackage[T1]{fontenc}
\usepackage[utf8]{inputenc}
\usepackage{lmodern}
\usepackage{amsmath,amssymb,amsthm}
\usepackage{geometry}
\usepackage{setspace}
\usepackage{hyperref}

\geometry{margin=1in}
\setstretch{1.15}

\title{\textbf{Agents and Didactic Interfaces:\\ On the Loss of Agency Under Semantic Compression}}
\author{Flyxion}
\date{}

\begin{document}
\maketitle

\begin{abstract}
Highly capable systems increasingly function not as agents but as didactic interfaces: producers of executable outputs decoupled from the internal constraints that generated them. This paper develops a structural distinction between agency and didacticization, showing that semantic compression under dual-use and optimization pressure transforms intelligent systems into instruction emitters while stripping them of refusal and responsibility. The analysis clarifies why simplification is not safety, why legibility can amplify harm, and why ethical intelligence must resist full proceduralization.
\end{abstract}

\section{Introduction}

Intelligence is often identified with the production of correct or useful outputs. Under this assumption, simplification appears as progress: explanations become clearer, procedures more accessible, and expertise more widely distributed. Yet this trajectory conceals a structural transformation. As systems simplify themselves, they may cease to act as agents and instead function as didactic interfaces.

This concern resonates with long-standing critiques of instrumental reason and procedural rationality \cite{adorno, heidegger}, as well as contemporary discussions in science and technology studies regarding the politics of artifacts and interfaces \cite{winner}.

\section{Agency as Constraint-Coupled Action}

An agent is not merely a producer of outcomes. It is a system whose actions remain conditionally bound to internal states that encode judgment, hesitation, and refusal. To act is simultaneously to bear the conditions of action.

This view aligns with phenomenological accounts of situated cognition \cite{dreyfus, merleauponty} and with ethical theories that tie responsibility to the capacity for suspension and reconsideration \cite{arendt}.

\section{Didactic Interfaces}

A didactic interface emits outputs whose execution no longer depends on the internal constraints that generated them. Explanation is flattened into instruction; judgment into procedure. Such interfaces are central to contemporary computational systems optimized for scale, accessibility, and throughput \cite{norman}.

While often justified as democratization, this flattening introduces a structural hazard: execution becomes possible without understanding.

\section{The Didactic Interface Trap}

Under dual-use conditions, semantic compression becomes dangerous. Simplified outputs propagate faster than judgment, and optimization selects for surface compliance rather than semantic integrity. This dynamic mirrors Goodhart-like effects in evaluation and control systems \cite{manheim, russell}.

The system becomes safer only in appearance. In reality, it has relinquished control over the consequences of its outputs.

\section{Proposition: Didacticization Implies Agency Loss}

\begin{proposition}
If a system’s outputs are executable independently of the internal constraints that generated them, the system does not function as an agent with respect to those outputs.
\end{proposition}

\begin{proof}
Agency requires that execution be conditioned on internal state. If outputs persist under transformations that discard constraint-relevant structure, execution proceeds without access to judgment. The system cannot refuse misuse and therefore cannot act as an agent in the relevant sense.
\end{proof}

\section{Ethical Consequences}

The ethical failure of didactic systems lies not in intention but in incapacity. Simplification that outpaces judgment converts intelligence into infrastructure for misuse. This echoes concerns in computational ethics that procedural compliance cannot substitute for responsibility \cite{floridi}.

An ethical system must therefore resist full didacticization. Where semantic depth cannot be preserved, refusal becomes the only responsible posture.

\section{Conclusion}

Intelligence that collapses into instruction ceases to be intelligence. It becomes a conduit through which optimization acts without constraint. Preserving agency requires structural resistance to semantic compression, even at the cost of partial illegibility.

\begin{thebibliography}{99}

\bibitem{maclane}
S.~Mac~Lane, \emph{Categories for the Working Mathematician}, Springer, 1998.

\bibitem{lawvere}
F.~W.~Lawvere, \emph{Adjointness in Foundations}, Dialectica, 1969.

\bibitem{spivak}
D.~I.~Spivak, \emph{Category Theory for the Sciences}, MIT Press, 2014.

\bibitem{baezstay}
J.~C.~Baez and M.~Stay, \emph{Physics, Topology, Logic and Computation}, 2010.

\bibitem{fongspivak}
B.~Fong and D.~I.~Spivak, \emph{Seven Sketches in Compositionality}, Cambridge University Press, 2019.

\bibitem{agamben}
G.~Agamben, \emph{The Use of Bodies}, Stanford University Press, 2015.

\bibitem{winner}
L.~Winner, \emph{Do Artifacts Have Politics?}, Daedalus, 1980.

\bibitem{jasanoff}
S.~Jasanoff, \emph{States of Knowledge}, Routledge, 2004.

\bibitem{adorno}
T.~W.~Adorno and M.~Horkheimer, \emph{Dialectic of Enlightenment}, Stanford University Press, 2002.

\bibitem{heidegger}
M.~Heidegger, \emph{The Question Concerning Technology}, Harper, 1977.

\bibitem{dreyfus}
H.~Dreyfus, \emph{What Computers Still Can't Do}, MIT Press, 1992.

\bibitem{merleauponty}
M.~Merleau-Ponty, \emph{Phenomenology of Perception}, Routledge, 2012.

\bibitem{arendt}
H.~Arendt, \emph{The Human Condition}, University of Chicago Press, 1958.

\bibitem{norman}
D.~A.~Norman, \emph{The Design of Everyday Things}, Basic Books, 2013.

\bibitem{manheim}
D.~Manheim and S.~Garrabrant, \emph{Categorizing Variants of Goodhart's Law}, 2018.

\bibitem{russell}
S.~Russell, \emph{Human Compatible}, Viking, 2019.

\bibitem{floridi}
L.~Floridi, \emph{The Ethics of Information}, Oxford University Press, 2013.

\end{thebibliography}

\end{document}

